\chapter{Trust-Aware GNN Inference and Task Assignment}
\label{ch:gnn}

\section{Role of the Inference Stage}

Pruning decides which links the fleet will use; inference decides which rovers
do the work. The stage runs on the fog device every cycle, so it must be cheap,
and it combines what a rover is (trust, battery, free compute, capability) with
where it sits in the kept graph (how reliable its neighbours are and how well
its results can reach the root).

\section{Node Features and Message Weights}

Each rover $i$ starts from the feature vector
\begin{equation}
h_i^{(0)}=\bigl[T_i,\ B_i,\ 1-\ell_i,\ q_i\bigr]\in[0,1]^4,
\end{equation}
trust, battery, available compute and capability. Messages along a kept edge
$(i,j)$ are weighted by
\begin{equation}
a_{ij}=T_{ij}\,L_{ij},
\end{equation}
the product of the edge trust (the smaller endpoint trust) and the measured
link quality. A neighbour reached over a high-trust, high-quality link
therefore contributes more to a rover's aggregate than one reached over a
poor link, and pruned edges contribute nothing at all, which is where the
pruner's decision enters the scores.

\section{Two-Layer Message Passing}

Each layer forms the trust-weighted mean of the neighbours' states and
combines it with the rover's own state:
\begin{equation}
m_i=\frac{\sum_{j}a_{ij}\,h_j}{\sum_j a_{ij}},\qquad
h_i'=\mathrm{ReLU}\bigl(W_sh_i+W_nm_i+b\bigr),
\label{eq:gnn:layer}
\end{equation}
with $m_i=0$ for a rover without kept neighbours. After two layers a
logistic readout produces the raw score
\begin{equation}
s_i=\sigma\bigl(w^{\top}h_i^{(2)}+\beta\bigr).
\label{eq:gnn:readout}
\end{equation}
The default weights are deliberately simple: $W_s=0.8\,I_4$,
$W_n=0.2\,I_4$, $b=0$, two layers, and a readout $w=(0.60,0.15,0.10,0.15)$
over trust, battery, compute and capability with $\beta=0$, dominated by
trust.



The self weight dominates for a reason found empirically: with equal
weighting, neighbour information washed out the per-rover differences that
assignment depends on, and every rover in a cluster scored alike.

\section{Gates, Criticality and Delivery Awareness}

A raw score is not yet a decision. Three adjustments turn it into one.

\subsubsection{Resource gates} A rover is ineligible for a task, with score
zero, if its trust is below the task's minimum (0.25 in the simulator), its
battery is below 15\% or its free compute is below the task's need (0.3).
Gates are hard because these failures are not matters of degree: a rover
with an empty battery cannot execute, however good its neighbourhood.

\subsubsection{Criticality} For mission-critical tasks (criticality above
0.5) the score is multiplied by $(0.5+0.5\,T_i)$. The scaling acts on the
probability rather than on the logit, so its direction does not depend on the
sign of the logit, and low-trust rovers are penalised more strongly exactly
when failure is most costly.

\subsubsection{Delivery awareness} The final score is multiplied by
$(0.5+0.5\,p_i)$, where $p_i$ is the delivery factor of
Chapter~\ref{ch:arch}: the weakest link on the widest path from $i$ to the
root in the kept graph. A rover that cannot reach the root at all ($p_i=0$)
is ineligible, because its result could never be delivered. The multiplier is
only partial so that the root, whose delivery factor is always one, does not
absorb every task merely by being the root.

\subsubsection{Deadlock fallback} If the trust gate would leave a task
without any eligible rover, for example early in a mission when every
estimate is still near the prior, the trust gate is relaxed for that task
while the resource gates are kept. The mission therefore never stalls on its
own caution.

\section{Capacity-Aware Greedy Assignment}

Given scores for every rover and task, the assignment is greedy and
conflict-free. All positive (score, task, rover) triples are sorted by
decreasing score, ties broken by task name for determinism, and each triple is
accepted if the task is still unassigned and the rover still has capacity
(one task per rover per cycle by default). Greedy matching is not optimal in
general, but with one task per rover and a handful of tasks per cycle it is
close, and it costs microseconds.

Each assignment raises the rover's load, which decays over the following
cycles; without this feedback the most trusted rover received every task until
its compute was exhausted.

\subsubsection{Exploration} With probability $\varepsilon=0.15$ per task,
the assigned rover is replaced by a random eligible one. Exploration keeps
trust estimates honest, because a rover that is never chosen is never
observed, and the explored assignments are tagged: they form the
random-policy, selection-bias-free subset used to train the readout in
Section~\ref{sec:gnn:readout}.

\section{The Online Trust Model}

Trust is learned from outcomes. After each cycle, every assigned rover is
updated with an exponentially weighted moving average,
\begin{equation}
T_i\leftarrow\alpha\,T_i+(1-\alpha)\,\mathrm{obs}_i,
\label{eq:gnn:ewma}
\end{equation}
where $\mathrm{obs}_i$ is 1 if the rover executed its task and 0 otherwise
(an optional link-evidence term can be mixed in with weights 0.7 and 0.3).
Rovers that were not observed age slowly towards the neutral prior,
\begin{equation}
T_i\leftarrow T_i+\delta\,(0.5-T_i),
\label{eq:gnn:aging}
\end{equation}
with $\alpha=0.5$ and $\delta=0.1$ in the simulator.

\begin{example}[Rehabilitation after one failure]
A rover at trust 0.5 fails once: Eq.~\eqref{eq:gnn:ewma} drops it to 0.25,
which sits on the trust gate. If it is then idle, aging restores it as
$0.5-0.25\times0.9^k$ after $k$ cycles, reaching 0.35 after five cycles and
0.41 after ten. Without aging the rover would sit at 0.25 forever, never be
chosen, and therefore never be observed again: a single unlucky failure would
blacklist it permanently. This was observed before aging was added.
\end{example}

The target the trust model tracks is deliberately hard. In the simulator,
hidden reliability is spatially correlated (nearby rovers share terrain and
interference) and can drift over time, so trust must follow a moving,
structured target rather than estimate a fixed coin bias. The spatial
correlation is also what gives message passing something to exploit:
neighbours' outcomes carry information about a rover that has not been
observed recently.

\section{Edge Lifecycle}

Pruned links are not forgotten. The candidate graph is rebuilt from
telemetry every cycle, so a pruned link can always return, and a small state
machine tracks each link's recent history: a link is \emph{active} while kept,
\emph{suspect} after one prune, in \emph{probation} after two and
\emph{removed} after four consecutive prunes, and returns to \emph{active} as
soon as it is kept again.



Links in \emph{probation} receive a freshness floor of 0.4. Without it, a
link that was pruned for a few cycles would also go unobserved, its freshness
would decay, its utility would fall further, and it would sink to the bottom
of the ranking for good. The floor gives such links a fair re-evaluation
when conditions improve.

\section{The Cloud-Trained Readout}
\label{sec:gnn:readout}

The readout $(w,\beta)$ of Eq.~\eqref{eq:gnn:readout} can be trained in the
cloud by logistic regression on pairs of hidden states and observed
outcomes, which would turn the hand-set scorer into a learned model while
keeping fog-side inference unchanged. Training data is subject to selection
bias, because outcomes are observed only for rovers the policy chose, so
training uses only the exploration-tagged samples, which were assigned at
random. With that unbiased but small dataset (78 samples), the trained
readout did \emph{not} improve held-out mission success over the fixed
weights (32.8\% against 36.7\%, a difference of $-3.9$ points with 95\% CI
$[-11.7,+11.7]$ and $p=0.49$). This is reported as a null result rather than
tuned until it looked better; a larger exploration budget is listed as future
work.

\section{Ablation Arms}

To isolate what each part contributes, the scorer can be replaced while gates,
delivery awareness and assignment stay identical: \emph{trust-only} (score
equals trust, the bandit-style baseline~\cite{auer2002}), \emph{no-graph} (the
same readout on raw features, without message passing) and \emph{random}.
Chapter~\ref{ch:results} reports their paired comparison, and the 3D engine
runs all three alongside a live mission.

\section{Computational Cost}

For $n$ rovers and $|E'|$ kept edges, one layer costs $O(n+|E'|)$ small
vector operations, and the whole stage measures about 0.2~ms at ten rovers
and about 1~ms at fifty on the evaluation machine. The constraint that
motivated the architecture is therefore met with a large margin: inference is
not the bottleneck of the pipeline, and Chapter~\ref{ch:results} shows that
pruning is.
